\scanpage{352}
% top centre
\mspage{268}
% lines 1–12
$C$ une catégorie, équivalente à une petite catégorie, $\pi$ un cardinal infini, posons $M=\operatorname{Ind}_\pi(C)$. Soit $\unsure{\Phi|\underline{\Phi}}$ un ens. de catégories de cardinal réduit $\leq\pi$. On se propose \deleted{\unsure{une famille}}\inserted{d'examiner} si $M$ est ``stable par $\varinjlim$ de type $\Phi$'' (i.e stable par $\unsure{\varinjlim|\underline{\lim}}$ de type $I$, quel que soit $I\in\Phi$). \uline{Il} \deleted{\unsure{Il|C}} \deleted{\unsure{en} est ainsi si} \inserted{suffit pour cela que} $C$ soit stable par $\varinjlim$\deleted{\unsure{ites}} de type $\Phi$, et que de plus : \deleted{$C$ est \unsure{st}} $\exists$ un ordinal \inserted{limite} $\alpha$ avec $\mathrm{card}\,\alpha\leq\pi$ et tel que $C$ soit stable par $\varinjlim$ de type $I_\alpha$ (p. ex. tel que $I_\alpha\in\Phi$).

% lines 13–20
\deleted{Cette condition} Inversement, si $M$ est stable par $\varinjlim$ de type $\Phi$, et si on suppose de plus $C$ \unsure{karoubienne} (i.e \unsure{stable|stables} \deleted{dans $M$} par facteurs directs) alors $C$ est stable par $\varinjlim$ de type $\Phi$. \deleted{Et le foncteur d'inclusion} Dans l'un et l'autre cas, \deleted{$C$ est}\inserted{le} foncteur d'inclusion $C\unsure{\lhook\joinrel\longrightarrow|\mathrel{-}}M$ commute aux $\varinjlim$ envisagées.
% lower left, diagonal beside conclusion
% slanted
\margin{\unsure{Explicité|Expliciter} \unsure{ds|cf} les corollaires \unsure{du|des} \unsure{type|types} 4.13.20\unsure{,|à|--} 4.13.21 (p. 341\unsure{)|}}

% lines 21–22
\uline{Démonstration} : \deleted{\unsure{Sous}} a) Condition suffisante. On applique 4.\unsure{90|9|10}.3 pour conclure que
% archive sheet number: 351
